% Generated by SIGMA; compile with Tectonic or another TeX engine.
\documentclass[11pt]{article}
\usepackage[T1]{fontenc}
\usepackage[utf8]{inputenc}
\usepackage[margin=1in]{geometry}
\usepackage{amsmath,amssymb,booktabs,longtable}
\usepackage{array}
\usepackage{enumitem}
\setlist{nosep,leftmargin=*}
\usepackage{microtype}
\usepackage{titlesec}
\usepackage{xcolor}
\definecolor{blue}{HTML}{24485E}
\titleformat{\section}{\Large\bfseries\sffamily\color{blue}}{\thesection}{0.8em}{}
\titleformat{\subsection}{\large\bfseries\sffamily}{\thesubsection}{0.7em}{}
\titleformat{\subsubsection}{\normalsize\bfseries\sffamily\color{blue}}{\thesubsubsection}{0.6em}{}
\setlength{\parskip}{0.55em}
\setlength{\parindent}{0pt}
\title{Maximum regions formed by straight lines}
\author{SIGMA Research Record}
\date{\today}
\begin{document}
\maketitle
\begin{quote}\small In the future we plan on using AI to produce a typical academic paper based on the results of the research.\end{quote}
\section{Project}
\textbf{Interpretation:} Agent-generated findings are provisional unless explicitly verified; reviewer agreement is not proof.\par
\textbf{Research problem:} For each nonnegative integer n, determine the greatest possible number of regions into which n distinct straight lines divide the Euclidean plane. Establish an exact answer for all n with a rigorous upper bound and a construction that attains it. State any geometric assumptions required for the argument.\par
\textbf{Project ID:} 9018b1fe-83a9-4027-a755-8f997b4f46b5\par
\textbf{Research type:} Mathematics\par
\textbf{Difficulty:} Easy\par
\section{Strategies}
\subsection{Incremental line insertion (saved)}
Examine how adding one line changes the connected regions of an arrangement of earlier lines. Identify what would need to be proved for an upper bound and what separate evidence would demonstrate attainability. Do not assume that a particular arrangement or formula is optimal.\par
\section{Research sessions}
\subsection{Session \#1}
\begin{description}[style=nextline,font=\normalfont\bfseries,leftmargin=0pt,labelsep=0pt]
\item[Execution status] ended
\item[Outcome] Mixed · unstructured argument, unassessed
\item[Strategy] Incremental line insertion
\item[Started] 2026-10-24T02:22:51.652165+00:00
\end{description}
\textbf{Outcome summary:} A solution candidate was recorded. Structured argument steps were not recorded; saved narrative is available for inspection in Reasoning / proof and the research artifacts. Its completeness and correctness remain unassessed.\par
\subsection{Conclusion}
\#\#\# Historical session

No Conclusion was saved for this session. Its existing outcome summary and detailed evidence remain available below.

Structured argument steps were not recorded; saved narrative is available for inspection in Reasoning / proof and the research artifacts. Its completeness and correctness remain unassessed.\par
Outcome caveats:\par
\begin{itemize}
\item Agent-reported evidence is unverified; agreement is not independent verification.
\item Historical or incomplete outcome record: no structured argument steps were preserved. Structured argument steps were not recorded; saved narrative is available for inspection in Reasoning / proof and the research artifacts. Its completeness and correctness remain unassessed. Agent agreement is not verification.
\end{itemize}
\textbf{Verification:} Convincing but unverified\par
Structured argument steps were not recorded; saved narrative is available for inspection in Reasoning / proof and the research artifacts. Its completeness and correctness remain unassessed.\par
\subsection{Outcome metrics}
\subsection{Operational metrics}
\textbf{Participating agents:} Investigator, Skeptic\par
\textbf{Model calls by status:} 4 dispatched · 4 completed · 0 rejected · 0 timed out · 0 interrupted / cancelled · 0 failed · 0 in progress · 0 not dispatched · 0 status unavailable\par
Tokens · input / output / cached / total: 32,042 / 6,549 / 0 / 38,591\par
Wall-clock / agent call time: 2m 29s wall · Investigator 45s across 2 call(s) · Skeptic 1m 42s across 2 call(s)\par
\textbf{Python runs by status:} 0 executions\par
\textbf{SageMath runs by status:} 0 executions\par
\textbf{Artifacts and storage:} 2 finding(s) · 0 Python file(s) (0 B) · session record 92,728 B\par
Provider failures / retries: 0 provider failure(s) (0 rejected, 0 timed out, 0 other) · 0 interrupted / cancelled · 0 later attempt(s); no automatic retries\par
\textbf{Estimated monetary cost:} — (plan allowance consumption unknown)\par
\subsection{Research-quality metrics}
\textbf{Strategies investigated:} Incremental line insertion\par
\textbf{Research cycles:} 2 completed / 2 started\par
\textbf{Outcome classification:} Mixed · unstructured argument, unassessed\par
\textbf{Verification status:} Convincing but unverified\par
\textbf{Open proof obligations:} 6 · Regions are open convex sign-pattern sets: established by the visible argument in cycle 1, artifact 2a734507-ad3f-4d5b-b00b-80b0e1624717, and repeated in cycle 2, step 9b40aecc-0b5f-40ff-aeeb-30faf598dcc4.; A distinct new line with t distinct old intersections creates exactly t+1 additional regions: established by the complete convexity and splitting argument in cycle 1, artifact 2a734507-ad3f-4d5b-b00b-80b0e1624717. (+4 more)\par
\textbf{Evidence-backed conclusions:} 2 conclusion(s) · 2 cited finding(s)\par
\textbf{Strategy contribution:} contributed to solution: 1\par
\textbf{Unresolved objections:} 0 · None recorded\par
\textbf{Claim scope:} not recorded\par
\textbf{Solution claim:} For every integer n≥0, the maximum number of connected regions formed by n distinct entire straight lines in the Euclidean plane is 1+n(n+1)/2. For n≥1, the lines y=kx+k², k=1,…,n, attain the maximum; the empty arrangement attains it for n=0.\par
\textbf{Argument:} Cycle 1 supplies the convex-region proof, the exact insertion increase t+1, and the universal upper bound. Cycle 2 supplies the missing exact intersection calculation and distinctness proof for the explicit construction. These establish equality at every insertion. The empty plane provides the base case. The complete conclusion is written out in this review without relying on any unseen continuation.\par
\textbf{Skeptic response:} The objective is achieved by combining the complete insertion argument and upper bound visible in cycle 1, artifact 2a734507-ad3f-4d5b-b00b-80b0e1624717, with the construction calculation supplied in cycle 2. The current output again ends before its final conclusion, but the missing conclusion follows directly from the supplied evidence; another cycle merely to restate it is unnecessary.

Here is the complete argument. Regions mean connected components of the complement of the union of the entire, pairwise distinct lines in the Euclidean plane. For an existing arrangement, its regions are precisely the nonempty fixed-sign sets of affine equations for its lines. Such sets are intersections of open half-planes, hence open and convex.

If a new line has t distinct intersection points with the previous lines, those points divide it into t+1 intervals, including the case t=0. Each interval belongs to one old region. Different intervals belong to different regions: otherwise convexity would force a segment containing an old intersection point to lie inside an old region. Conversely, every old region meeting the new line contains one of these intervals. Thus exactly t+1 old regions meet the new line. Each is split into exactly two nonempty convex pieces: openness supplies points on both sides, and intersection with either open half-plane is convex. Other regions are unchanged, and deleting additional points from the complement cannot merge old components. Consequently the increase is exactly t+1.

At insertion k there are at most k−1 distinct intersections with earlier lines, so the increase is at most k. Starting with the single region of the empty arrangement gives the universal upper bound 1+Σ(k=1 to n)k = 1+n(n+1)/2. Neither parallel pairs nor multiple concurrence invalidate this bound.

For attainment, take L\_k: y=kx+k² for k=1,…,n. For i\textless{}k, solving the two equations gives L\_i∩L\_k=\{(-(i+k),-ik)\}. For fixed k, distinct i give distinct x-coordinates. Therefore L\_k has exactly k−1 distinct intersections with its predecessors and adds exactly k regions. This construction consequently has 1+n(n+1)/2 regions. For n=0, the complement is the whole plane and has one region, agreeing with the same formula.

Hence the greatest possible number is M(n)=1+n(n+1)/2 for every nonnegative integer n. Distinct slopes and absence of triple concurrence are properties of the attaining construction, not extra assumptions on arbitrary arrangements. No computational evidence is used\par
\textbf{Assumptions reviewed:} The ambient space is the ordinary Euclidean plane.; Each line is an entire infinite straight line, and the n lines are pairwise distinct.; A region is a connected component of the complement of the union of the lines.; The upper bound allows parallel pairs and multiple concurrence.; The proposed construction has distinct slopes and distinct predecessor-intersection points on each inserted line; these properties have been justified algebraically.\par
\textbf{Proof obligations considered:} Regions are open convex sign-pattern sets: established by the visible argument in cycle 1, artifact 2a734507-ad3f-4d5b-b00b-80b0e1624717, and repeated in cycle 2, step 9b40aecc-0b5f-40ff-aeeb-30faf598dcc4.; A distinct new line with t distinct old intersections creates exactly t+1 additional regions: established by the complete convexity and splitting argument in cycle 1, artifact 2a734507-ad3f-4d5b-b00b-80b0e1624717.; Every arrangement satisfies R≤1+n(n+1)/2: established by the insertion count and summation in cycle 1, artifact 2a734507-ad3f-4d5b-b00b-80b0e1624717.; An attaining arrangement exists for every positive n: cycle 2, step 9b40aecc-0b5f-40ff-aeeb-30faf598dcc4, proves L\_i∩L\_k=\{(-(i+k),-ik)\} and distinctness for fixed k. Together with the cycle 1 insertion lemma, this establishes equality at every insertion.; The formula holds for n=0 and is the exact maximum for all n: cycle 1 explicitly supplies the one-region empty-arrangement base case and upper bound; cycle 2 supplies attainment. Their combination, explicitly concluded in this review, establishes the result.\par
\textbf{Still open:} Finish and state the insertion lemma, sum the upper bound, and explicitly conclude attainment and n = 0.\par
\textbf{Solution evidence references:} 2a734507-ad3f-4d5b-b00b-80b0e1624717; a832588d-9da7-4ceb-b3ff-9451a9922ff5\par
\textbf{Unresolved objections:} None recorded\par
\textbf{Resolved review notes:} Material and recoverable: the supplied construction paragraph is incomplete. Although a promising family is displayed, the visible argument does not establish that its intersections with each new line are distinct, so attainability has not yet been demonstrated in the supplied proof.; The visible proof does not finish checking the proposed attaining construction.; Complete the attainment proof for L\_k: y = kx + k² by checking pairwise intersections and distinctness of the intersection points on each inserted line.; Combine attainment with the established upper bound and explicitly cover n = 0.\par
\textbf{Saved verification code:} convincing\_unverified\par
\textbf{Strategy finding:} Incremental line insertion · achieved\par
\textbf{Strategy execution status:} ended\par
\textbf{Strategy failure detail:} Stop productive conclusion\par
\textbf{Claim:} For every integer n≥0, the maximum number of connected regions formed by n distinct entire straight lines in the Euclidean plane is 1+n(n+1)/2. For n≥1, the lines y=kx+k², k=1,…,n, attain the maximum; the empty arrangement attains it for n=0.\par
\textbf{Source artifacts:} 2a734507-ad3f-4d5b-b00b-80b0e1624717; a832588d-9da7-4ceb-b3ff-9451a9922ff5\par
\textbf{Open obligation:} Finish and state the insertion lemma, sum the upper bound, and explicitly conclude attainment and n = 0.\par
\subsubsection*{Strategy assessments}
\begin{longtable}{>{\raggedright\arraybackslash}p{.16\linewidth}>{\raggedright\arraybackslash}p{.15\linewidth}>{\raggedright\arraybackslash}p{.40\linewidth}>{\raggedright\arraybackslash}p{.17\linewidth}}
\toprule Strategy & Assessment & Rationale & Evidence \\
\midrule\endhead
Incremental line insertion & Contributed to solution & Incremental line insertion supplies both a universal upper bound and an exact equality criterion. The explicit family satisfies that criterion for every insertion, completing the selected strategy without numerical testing or an additional geometric existence assumption. & 2a734507-ad3f-4d5b-b00b-80b0e1624717; a832588d-9da7-4ceb-b3ff-9451a9922ff5 \\
\bottomrule
\end{longtable}
\paragraph{Artifact 1 · Unverified research note. Any Investigator Python or SageMath execution, including its source, status, output, and captured artifacts, is recorded below and is not proof by itself.}
\textbf{Claim:} The greatest number of regions is M(n) = 1 + n(n+1)/2 for every integer n ≥ 0. The key incremental lemma is exact: if a new line meets the union of the earlier lines in t distinct points, it creates precisely t+1 additional regions.\par
\textbf{Method:} Hypothesis and discriminating criterion: the kth inserted line creates at most k regions; equality at every insertion would establish attainability of the resulting upper bound. The exact insertion lemma below proves this criterion without a numerical experiment.

1. Convexity of regions. Choose an affine linear equation for each existing line. Every point off the lines has a sign for each equation. Points with a fixed sign pattern form an intersection of open half-planes, hence a convex set when nonempty. Such a set is connected, and signs cannot change along a connected subset of the complement. Therefore the regions are exactly the nonempty sign-pattern sets and are open and convex.

2. Exact insertion lemma. Let L be a line distinct from all existing lines. Its t distinct intersection points with their union divide L into t+1 open intervals, including the unbounded intervals. Each interval lies in a single old region. Different intervals lie in different old regions: otherwise convexity would force the segment joining points in those intervals to remain in that region, although the segment contains an old line intersection. Conversely, any old region meeting L contains a point of one of these intervals, since it contains no point on an old line. Thus exactly t+1 old regions meet L. Each such region is split into exactly two regions: its intersections with the two open half-planes bounded by L are convex, and both are nonempty because the old region is open and contains a point of L. Regions not meeting L remain unchanged. The increase is consequently exactly t+1.

3. Upper bound. At insertion k, the new line has at most k−1 distinct intersection points with earlier lines, because distinct Euclidean lines meet in at most one point. It therefore adds at most k regions. Starting with one region for no lines gives R ≤ 1 + Σ(k=1 to n) k = 1 + n(n+1)/2.

4. Attaining construction. For n ≥ 1 take L\_k: y = kx + k², for k = 1,…,n, and insert them in this order. Their pen\par
\textbf{Assumptions:} The ambient space is the ordinary Euclidean plane.; Each line is an entire infinite straight line, and the n lines are pairwise distinct.; A region means a connected component of the complement of the union of the lines.; Parallel lines and multiple concurrence are permitted for the upper bound; the attaining construction has neither parallel pairs nor triple concurrence.\par
\textbf{Evidence status:} A deductive proof is supplied, rather than a conjecture or finite computational test. Its essential geometric step is justified by convexity of the complement regions. No external checking or formal verification is claimed.\par
\textbf{Skeptic response:} The supplied argument establishes the incremental lemma and the upper bound M(n) ≤ 1 + n(n+1)/2. The region-convexity argument correctly handles arbitrary arrangements, including parallel lines and multiple concurrence. However, the visible attainment paragraph ends mid-sentence after proposing L\_k: y = kx + k². I cannot treat its unseen continuation as evidence that the construction was justified. Please complete the short exact check that each inserted L\_k meets its k−1 predecessors at distinct points, and then explicitly conclude equality for all n, including n = 0. This is a recoverable proof-completion issue, not a counterexample or a request for computation.\par
\textbf{Objections:} Material and recoverable: the supplied construction paragraph is incomplete. Although a promising family is displayed, the visible argument does not establish that its intersections with each new line are distinct, so attainability has not yet been demonstrated in the supplied proof.\par
\textbf{Decision:} continue\par
\textbf{Open obligations:} Complete the attainment proof for L\_k: y = kx + k² by checking pairwise intersections and distinctness of the intersection points on each inserted line.; Combine attainment with the established upper bound and explicitly cover n = 0.\par
\textbf{Next question:} For the proposed lines L\_k: y = kx + k², compute L\_i ∩ L\_k for i \textless{} k and prove that these k−1 intersection points are distinct; then finish the equality argument, including n = 0.\par
\paragraph{Artifact 2 · Unverified research note. Any Investigator Python or SageMath execution, including its source, status, output, and captured artifacts, is recorded below and is not proof by itself.}
\textbf{Claim:} For every nonnegative integer n, the greatest possible number of regions is M(n) = 1 + n(n+1)/2. For n ≥ 1, the lines L\_k: y = kx + k², k = 1,…,n, attain this bound.\par
\textbf{Method:} Criterion: the kth inserted line adds at most k regions. To prove attainment, it suffices to show that each proposed L\_k meets its k−1 predecessors at distinct points. The following exact argument establishes both statements; no numerical experiment is needed.

1. Intersection calculation and distinctness. For 1 ≤ i \textless{} k, a point on both L\_i and L\_k satisfies ix + i² = kx + k². Thus (i−k)x = (k−i)(k+i), and division by i−k ≠ 0 gives x = −(i+k). Substituting into either equation gives y = −ik. Therefore L\_i ∩ L\_k = \{ (−(i+k), −ik) \}. For fixed k, distinct indices i give distinct x-coordinates, so these k−1 intersection points on L\_k are all distinct. The slopes 1,…,n are distinct as well, so the lines are distinct and no pair is parallel. There is no triple concurrence: the largest-index line of any purported concurrent triple would have two identical intersection points with its predecessors, contradicting the distinctness just proved.

2. Exact insertion lemma. Regions of a finite line arrangement are open convex sets. Indeed, choose an affine linear equation for each line. Each nonempty fixed-sign set is an intersection of open half-planes and hence is open and convex. It is connected, while every connected subset of the complement has constant signs. Thus these sets are precisely the connected components of the complement.

Suppose a new, distinct line L has t distinct intersection points with the earlier lines. These points cut L into t+1 open intervals, with the whole line treated as one interval if t = 0. Each interval lies in an old region. Two different intervals cannot lie in the same old region: convexity would put the segment joining points in them inside that region, although this segment contains an intersection with an old line. Conversely, every old region meeting L contains a point in one of these intervals. Exactly t+1 old regions therefore meet L. Each is split into exactly two nonempty convex pieces by L: both sides are nonempty because the old in­\par
\textbf{Assumptions:} The ambient space is the ordinary Euclidean plane.; The lines are entire infinite straight lines and are pairwise distinct.; A region is a connected component of the complement of the union of the lines.; The universal upper bound permits parallel pairs and multiple concurrence; no general-position assumption is required for that bound.\par
\textbf{Evidence status:} The intersection and distinctness calculations above are deductive. They complete the previously unfinished construction check in the supplied extract of artifact 2a734507-ad3f-4d5b-b00b-80b0e1624717. To finish the insertion argument: each old region meeting L is open and contains a point of L, so it has points on both sides of L. Its two intersections with the open half-planes bounded by L are convex and connected. Different old regions cannot merge when another line is removed from their union\par
\textbf{Skeptic response:} The objective is achieved by combining the complete insertion argument and upper bound visible in cycle 1, artifact 2a734507-ad3f-4d5b-b00b-80b0e1624717, with the construction calculation supplied in cycle 2. The current output again ends before its final conclusion, but the missing conclusion follows directly from the supplied evidence; another cycle merely to restate it is unnecessary.

Here is the complete argument. Regions mean connected components of the complement of the union of the entire, pairwise distinct lines in the Euclidean plane. For an existing arrangement, its regions are precisely the nonempty fixed-sign sets of affine equations for its lines. Such sets are intersections of open half-planes, hence open and convex.

If a new line has t distinct intersection points with the previous lines, those points divide it into t+1 intervals, including the case t=0. Each interval belongs to one old region. Different intervals belong to different regions: otherwise convexity would force a segment containing an old intersection point to lie inside an old region. Conversely, every old region meeting the new line contains one of these intervals. Thus exactly t+1 old regions meet the new line. Each is split into exactly two nonempty convex pieces: openness supplies points on both sides, and intersection with either open half-plane is convex. Other regions are unchanged, and deleting additional points from the complement cannot merge old components. Consequently the increase is exactly t+1.

At insertion k there are at most k−1 distinct intersections with earlier lines, so the increase is at most k. Starting with the single region of the empty arrangement gives the universal upper bound 1+Σ(k=1 to n)k = 1+n(n+1)/2. Neither parallel pairs nor multiple concurrence invalidate this bound.

For attainment, take L\_k: y=kx+k² for k=1,…,n. For i\textless{}k, solving the two equations gives L\_i∩L\_k=\{(-(i+k),-ik)\}. For fixed k, distinct i give distinct x-coordinates. Therefore L\_k has exactly k−1 distinct intersections with its predecessors and adds exactly k regions. This construction consequently has 1+n(n+1)/2 regions. For n=0, the complement is the whole plane and has one region, agreeing with the same formula.

Hence the greatest possible number is M(n)=1+n(n+1)/2 for every nonnegative integer n. Distinct slopes and absence of triple concurrence are properties of the attaining construction, not extra assumptions on arbitrary arrangements. No computational evidence is used\par
\textbf{Objections:} None recorded\par
\textbf{Decision:} stop\_productive\_conclusion\par
\textbf{Open obligations:} Finish and state the insertion lemma, sum the upper bound, and explicitly conclude attainment and n = 0.\par
\textbf{Next question:} Not recorded\par
\textbf{Obligation disposition:} addressed · 727f5797506a4ed9d396013a2a5e9fb53422493fbda677e1680ec89e19afa392 · Cycle 2 explicitly computes the intersection as (-(i+k),-ik) and proves distinctness by the x-coordinate for fixed k. This is a model-reviewed deductive resolution, not independent formal verification. · evidence: 9b40aecc-0b5f-40ff-aeeb-30faf598dcc4 · model-reported, unverified\par
\textbf{Obligation disposition:} addressed · 929448a83f82f9c479fca1a6574df242ac58dad95706df537e02446bc3cd08d5 · Cycle 1 supplies the upper bound and the empty-arrangement base case; cycle 2 supplies the construction check. Their equality conclusion is explicitly completed in this review. · evidence: 2a734507-ad3f-4d5b-b00b-80b0e1624717; 9b40aecc-0b5f-40ff-aeeb-30faf598dcc4 · model-reported, unverified\par
\textbf{Obligation disposition:} addressed · 0d5ce1980c1b9045ed9fc42c83eb9f6bf417cfbf7f3728d582af12e195f41c89 · The previously missing construction calculation is now visible and correct. The original incompleteness objection was recoverable, not a counterexample. · evidence: 9b40aecc-0b5f-40ff-aeeb-30faf598dcc4 · model-reported, unverified\par
\textbf{Obligation disposition:} addressed · efbceb9d366c84a5d396fbb42c1c3a2c33a63b18c6cb0f3feb817e6f555104b6 · Cycle 2 finishes checking the proposed construction by establishing distinct slopes and distinct intersections with predecessors. · evidence: 9b40aecc-0b5f-40ff-aeeb-30faf598dcc4 · model-reported, unverified\par
\textbf{Obligation disposition:} addressed · 0e3dacc0f6887835ec947e229ec09c2a01d11a561edda4be62996a938dece992 · The cycle 1 extract contains the actual sign-pattern proof of openness, convexity, and identification with connected components; this judgment does not rely solely on the earlier review. · evidence: 2a734507-ad3f-4d5b-b00b-80b0e1624717 · model-reported, unverified\par
\textbf{Obligation disposition:} addressed · c76dc03a8d8cd495ea8b865ad4955306785f815b88bf0b64831a1ccadf02bed7 · The complete cycle 1 argument establishes the interval-to-region correspondence and the split into two nonempty convex components, yielding exactly t+1 additional regions. · evidence: 2a734507-ad3f-4d5b-b00b-80b0e1624717 · model-reported, unverified\par
\textbf{Obligation disposition:} addressed · d7b59964436468abf1567b21b118dfdf93704a892084f01e37d631f2d78786d2 · Cycle 1 explicitly bounds the kth increase by k and sums from the one-region empty arrangement. · evidence: 2a734507-ad3f-4d5b-b00b-80b0e1624717 · model-reported, unverified\par
\textbf{Obligation disposition:} addressed · f7e7db1e14b3bccf940a1b0675c0877fd3f1cfd451f5d306102b8ac4ae3456b1 · The cycle 2 construction calculation combined with the cycle 1 exact insertion lemma establishes attainment for every positive n, with the empty arrangement covering zero. · evidence: 2a734507-ad3f-4d5b-b00b-80b0e1624717; 9b40aecc-0b5f-40ff-aeeb-30faf598dcc4 · model-reported, unverified\par
\textbf{Obligation disposition:} addressed · 3915952af2ba8049404ed909f2bb5879518cf0c38fabc08a02da3fc1cc1adc18 · The supplied arguments jointly establish the exact answer, explicitly concluded for all nonnegative n in this review. · evidence: 2a734507-ad3f-4d5b-b00b-80b0e1624717; 9b40aecc-0b5f-40ff-aeeb-30faf598dcc4 · model-reported, unverified\par
\textbf{Obligation disposition:} addressed · 17b7e633ca2544b15fc1d4bbdf224d53126053294b5f905d8ac9103e674b2e3d · The source pointer alone is not used as evidence. The separately supplied cycle 1 proof already finishes the insertion lemma and summation, while the actual cycle 2 output proves the construction criterion. This review explicitly combines those supplied ingredients and covers n=0; no unseen continuation is assumed. · evidence: 2a734507-ad3f-4d5b-b00b-80b0e1624717; 9b40aecc-0b5f-40ff-aeeb-30faf598dcc4 · model-reported, unverified\par
\subsection{Agent context selection}
Context extracts are deterministic, untrusted summaries; omitted evidence remains in the original session records.\par
\textbf{Step d9aeefa9-1c04-4615-883f-f58c4d83df9d · investigator:} 0 selected, 0 current-session artifacts omitted, 5100 UTF-8 bytes\par
\textbf{Obligation ledger dispositions:} unresolved 0, addressed 0, superseded 0, retracted 0, unknown 0\par
\begin{itemize}
\item selected project\_strategy\_snapshot: ["8dda06ce-c931-42a8-a2b8-a68f1220e848"] · bounded portfolio overview for strategy and domain context
\end{itemize}
\textbf{Step 5b7e08aa-c6aa-444a-8db9-b79fbdd9b919 · skeptic:} 0 selected, 0 current-session artifacts omitted, 8281 UTF-8 bytes\par
\textbf{Obligation ledger dispositions:} unresolved 0, addressed 0, superseded 0, retracted 0, unknown 0\par
\begin{itemize}
\item selected project\_strategy\_snapshot: ["8dda06ce-c931-42a8-a2b8-a68f1220e848"] · bounded portfolio overview for strategy and domain context
\end{itemize}
\textbf{Step 9b40aecc-0b5f-40ff-aeeb-30faf598dcc4 · investigator:} 1 selected, 0 current-session artifacts omitted, 35164 UTF-8 bytes\par
\textbf{Obligation ledger dispositions:} unresolved 9, addressed 0, superseded 0, retracted 0, unknown 0\par
\begin{itemize}
\item selected current\_session\_material\_evidence: ["2a734507-ad3f-4d5b-b00b-80b0e1624717"] · representative sources for unresolved obligations, structured decisive outcomes, and newest saved artifact; selected before optional history
\item selected project\_strategy\_snapshot: ["8dda06ce-c931-42a8-a2b8-a68f1220e848"] · bounded portfolio overview for strategy and domain context
\end{itemize}
\textbf{Step b604a89d-6ebc-4059-a4c9-64b095c27605 · skeptic:} 1 selected, 0 current-session artifacts omitted, 40095 UTF-8 bytes\par
\textbf{Obligation ledger dispositions:} unresolved 10, addressed 0, superseded 0, retracted 0, unknown 0\par
\begin{itemize}
\item selected current\_session\_material\_evidence: ["2a734507-ad3f-4d5b-b00b-80b0e1624717"] · representative sources for unresolved obligations, structured decisive outcomes, and newest saved artifact; selected before optional history
\item selected project\_strategy\_snapshot: ["8dda06ce-c931-42a8-a2b8-a68f1220e848"] · bounded portfolio overview for strategy and domain context
\item selected current\_investigator\_obligation: ["9b40aecc-0b5f-40ff-aeeb-30faf598dcc4"] · current Investigator step is the source; disposition remains unverified
\end{itemize}
\textbf{Terminal reason:} stop\_productive\_conclusion\par
\section{Agentic runs}
No agentic runs are recorded.\par
\end{document}
